\section{Methods}\label{sec:methods}

\paragraph{Learning problem.} Let $x_t\in\mathbb{R}^{57}$ be the features known at the close of day $t$.
The main task is to estimate $\pi(x_t)=\Pr(S_t=1\mid x_t)$, the probability of a spike; the secondary
task estimates the conditional mean premium $m(x_t)=\E[R_t\mid x_t]$. Both are supervised problems
solved by empirical risk minimisation, $\hat f=\arg\min_{f\in\mathcal F}\frac1n\sum_t L(y_t,f(x_t))+\lambda\,\Omega(f)$,
with the log loss $L=-[y\log f+(1-y)\log(1-f)]$ for the spike and the squared loss for the premium.
The log loss is a \emph{strictly proper} scoring rule \citep{gneiting2007}: it is minimised in expectation
only by the true probability, which is what the hedging rule needs (Appendix~\ref{app:theory}).
The three model families below span the bias--variance trade-off: from a rigid linear model to flexible
tree ensembles. Each is tuned on its own training window only (Section~\ref{sec:evaluation}).

\paragraph{Regularised logistic and ridge regression (baseline).} The logistic model assumes the
log-odds are \emph{linear and additive} in the standardised features,
$\log\frac{\pi(x)}{1-\pi(x)}=\beta_0+x^\top\beta$, with an $L_2$ penalty $\frac{1}{2C}\lVert\beta\rVert^2$; ridge
regression does the same for the premium with penalty $\alpha\lVert\beta\rVert^2$. The penalty shrinks
coefficients towards zero, trading a little bias for a large reduction in variance, which matters here
because many features are strongly correlated (returns over overlapping windows, several volatility
measures). Assumptions: linearity in the log-odds, no interactions unless engineered, features on a common
scale (hence standardisation, fitted on the training window only). It is the natural ``simple model'' a
risk desk would build and the reference for the more flexible models.

\paragraph{Random forest.} A forest averages $B=300$ classification (or regression) trees, each grown on a
random half of the training rows with a random subset of $\sqrt{p}$ features considered at each split
\citep{breiman2001}. A single tree partitions the feature space into rectangles and predicts the average
outcome in each; it captures non-linearities and interactions (e.g.\ \emph{high volatility and backwardation})
but has high variance. Averaging $B$ trees with pairwise correlation $\rho$ reduces the variance of the
ensemble to $\rho\sigma^2+\frac{1-\rho}{B}\sigma^2$; random feature selection lowers $\rho$. Assumptions:
none on functional form, but the forest cannot extrapolate beyond the range of the training data (a problem
in 2022, see Section~\ref{sec:limits}). Tuned: maximum depth $\{3,6\}$ and minimum leaf size $\{25,100\}$,
which control how finely the forest can fit the overlapping, autocorrelated daily labels.

\paragraph{Gradient boosting.} Boosting builds an additive model $F_M(x)=\sum_{m=1}^M\nu\,h_m(x)$ of
\emph{shallow} trees, each fitted to the negative gradient of the loss of the current model
\citep{friedman2001}; for the log loss the gradient is simply $y-\hat\pi$, so each tree corrects the
current probability errors. Where the forest reduces variance by averaging, boosting reduces bias step by
step; it is controlled by the learning rate $\nu=0.05$, an $L_2$ penalty on leaf values, the number of trees
$M\in\{100,300\}$ and the depth $\{2,3\}$ (depth $d$ allows interactions of order $d$). We use the
histogram-based implementation of scikit-learn, which bins features as in LightGBM \citep{ke2017} and
handles missing values natively.

\paragraph{Why these choices.} The sample has 4\,700 daily rows but, because labels overlap over 21 days,
only about 220 independent months. Grids are therefore deliberately small (4 configurations per model),
trees shallow, and all models regularised. We did not use neural networks: with this effective sample size
their extra flexibility would mostly add variance \citep{gu2020}.
